% GENERATED FILE. Edit canonical lessons, then run npm run generate:books.
% canonical-source-hash: fnv1a64:65f1b6c13c088ae5
% canonical-lessons: TA-C137-motiram, TA-C137-valaiyal, TA-C137-kalurai, TA-C137-kaiyurai, TA-C137-vakuppu

\chapter{Ring, Bangle, Sock, Glove, Class}
\label{ch:ta-ring-bangle-sock-glove-class}
\hlchaptermodality{137}

\begin{chapteropening}
\textbf{By the end of this chapter:} \emph{I can say a ring, a bangle, a sock, a glove and a class.}

Complete the last lesson of chapter 137: i can say a ring, a bangle, a sock, a glove and a class.
\end{chapteropening}

\section[mōtiram]{\ta{மோதிரம்} (mōtiram) — a ring}

\label{lesson:TA-C137-motiram}



\begin{quote}
Before the new one: say the Tamil for a doll, then the Tamil for a purse.
\end{quote}

\subsection*{You'll want to know: \ta{மோதிரம்}}
\textbf{\ta{மோதிரம்}} — \emph{mōtiram} — \textquotedblleft{}a ring\textquotedblright{}.

\begin{grammarlens}[title={What you've built}]
One more word for this chapter.
\end{grammarlens}

\subsection*{Your turn}
\begin{itemize}
\raggedright
  \item \emph{Say it:} \emph{mōtiram}
  \item \emph{Say it:} \emph{mōtiram}, once more
  \item \emph{Say it:} say \emph{mōtiram}
  \item \emph{Recall:} say the Tamil for a doll, then the Tamil for a purse, then say \emph{mōtiram} again
\end{itemize}

\begin{tcolorbox}[breakable,colback=teal!4,colframe=teal!35!black,title={Before you move on}]
What does \ta{மோதிரம்} mean? (\textquotedblleft{}A ring\textquotedblright{}.) Say it once more.
\end{tcolorbox}
\section[vaḷaiyal]{\ta{வளையல்} (vaḷaiyal) — a bangle}

\label{lesson:TA-C137-valaiyal}



\begin{quote}
Before the new one: say the Tamil for a purse, then the Tamil for a ring.
\end{quote}

\subsection*{You'll want to know: \ta{வளையல்}}
\textbf{\ta{வளையல்}} — \emph{vaḷaiyal} — \textquotedblleft{}a bangle\textquotedblright{}.

\begin{grammarlens}[title={What you've built}]
One more word for this chapter.
\end{grammarlens}

\subsection*{Your turn}
\begin{itemize}
\raggedright
  \item \emph{Say it:} \emph{vaḷaiyal}
  \item \emph{Say it:} \emph{vaḷaiyal}, once more
  \item \emph{Say it:} say \emph{vaḷaiyal}
  \item \emph{Recall:} say the Tamil for a purse, then the Tamil for a ring, then say \emph{vaḷaiyal} again
\end{itemize}

\begin{tcolorbox}[breakable,colback=teal!4,colframe=teal!35!black,title={Before you move on}]
What does \ta{வளையல்} mean? (\textquotedblleft{}A bangle\textquotedblright{}.) Say it once more.
\end{tcolorbox}
\section[kāluṟai]{\ta{காலுறை} (kāluṟai) — a sock}

\label{lesson:TA-C137-kalurai}



\begin{quote}
Before the new one: say the Tamil for a ring, then the Tamil for a bangle.
\end{quote}

\subsection*{You'll want to know: \ta{காலுறை}}
\textbf{\ta{காலுறை}} — \emph{kāluṟai} — \textquotedblleft{}a sock\textquotedblright{}.

\begin{grammarlens}[title={What you've built}]
One more word for this chapter.
\end{grammarlens}

\subsection*{Your turn}
\begin{itemize}
\raggedright
  \item \emph{Say it:} \emph{kāluṟai}
  \item \emph{Say it:} \emph{kāluṟai}, once more
  \item \emph{Say it:} say \emph{kāluṟai}
  \item \emph{Recall:} say the Tamil for a ring, then the Tamil for a bangle, then say \emph{kāluṟai} again
\end{itemize}

\begin{tcolorbox}[breakable,colback=teal!4,colframe=teal!35!black,title={Before you move on}]
What does \ta{காலுறை} mean? (\textquotedblleft{}A sock\textquotedblright{}.) Say it once more.
\end{tcolorbox}
\section[kaiyuṟai]{\ta{கையுறை} (kaiyuṟai) — a glove}

\label{lesson:TA-C137-kaiyurai}



\begin{quote}
Before the new one: say the Tamil for a bangle, then the Tamil for a sock.
\end{quote}

\subsection*{You'll want to know: \ta{கையுறை}}
\textbf{\ta{கையுறை}} — \emph{kaiyuṟai} — \textquotedblleft{}a glove\textquotedblright{}.

\begin{grammarlens}[title={What you've built}]
One more word for this chapter.
\end{grammarlens}

\subsection*{Your turn}
\begin{itemize}
\raggedright
  \item \emph{Say it:} \emph{kaiyuṟai}
  \item \emph{Say it:} \emph{kaiyuṟai}, once more
  \item \emph{Say it:} say \emph{kaiyuṟai}
  \item \emph{Recall:} say the Tamil for a bangle, then the Tamil for a sock, then say \emph{kaiyuṟai} again
\end{itemize}

\begin{tcolorbox}[breakable,colback=teal!4,colframe=teal!35!black,title={Before you move on}]
What does \ta{கையுறை} mean? (\textquotedblleft{}A glove\textquotedblright{}.) Say it once more.
\end{tcolorbox}
\section[vakuppu]{\ta{வகுப்பு} (vakuppu) — a class}

\label{lesson:TA-C137-vakuppu}



\begin{quote}
Before the new one: say the Tamil for a sock, then the Tamil for a glove.

Say \textbf{\ta{திங்கள்}} through \textbf{\ta{ஞாயிறு}}. Each day ends in \emph{kiḻamai}; some days carry Tamil words, some Sanskrit ones.
\end{quote}

\subsection*{You'll want to know: \ta{வகுப்பு}}
\textbf{\ta{வகுப்பு}} — \emph{vakuppu} — \textquotedblleft{}a class\textquotedblright{}.

\begin{grammarlens}[title={What you've built}]
One more word for this chapter.
\end{grammarlens}

\subsection*{Your turn}
\begin{itemize}
\raggedright
  \item \emph{Say it:} \emph{vakuppu}
  \item \emph{Say it:} \emph{vakuppu}, once more
  \item \emph{Say it:} say \emph{vakuppu}
  \item \emph{Recall:} say the Tamil for a sock, then the Tamil for a glove, then say \emph{vakuppu} again
\end{itemize}

\begin{tcolorbox}[breakable,colback=teal!4,colframe=teal!35!black,title={Before you move on}]
What does \ta{வகுப்பு} mean? (\textquotedblleft{}A class\textquotedblright{}.) Say it once more.
\end{tcolorbox}
